\section{Limitations and Future Work}

% In this section, we discuss limitations of our work and potential avenues for future work.

\subsection{Limitations}

\textbf{Limited Language Coverage.} Although \benchmarkname~includes multiple programming languages (Python, JavaScript, TypeScript, Go), the distribution is not uniform, and some widely-used languages like Java, C++, and Rust are underrepresented. This may limit the benchmark's ability to assess agent performance across the full spectrum of modern software development.

% \textbf{Issue Scope.} The current evaluation framework focuses primarily on issue resolution through code patches. Real-world software engineering encompasses broader activities such as system design, code review, documentation, and architectural decisions that are not captured in the current benchmark structure.

\textbf{Dependency on Test Suite.} We rely on a test suite of \texttt{fail2pass} and \texttt{pass2pass} to verify problem solutions. However, real software engineering tasks may have a variety of correct solutions, even if they do not pass the original tests outlined in the task. Ideally, we might have a set of verifiers which can verify any valid solution. 

% \textbf{Reduction in Ambiguity.} The human augmentation process, while improving problem clarity, may inadvertently make problems too prescriptive by providing excessive detail in requirements and interface specifications. In the real-world, problems are ambiguous, with potential follow-up or exploration needed to start the task.

\subsection{Future Work}

% \textbf{Expanded Language Coverage.} Future iterations of \benchmarkname~should incorporate more diverse programming languages and frameworks to better represent the software development ecosystem. This includes languages like Java, C\#, Rust, Kotlin, and emerging languages that may become prevalent in industry settings.

\textbf{Alternative Evaluation Metrics.} Developing evaluation approaches beyond test-based verification, such as rubrics, code quality assessment, security analysis, performance optimization, and adherence to software engineering best practices. This could include human evaluation of code maintainability, readability, and architectural soundness.

\textbf{Collaborative Development Scenarios.} Introducing problems that require coordination between multiple agents or human-agent collaboration, reflecting modern team-based software development practices. This could include scenarios involving code reviews, merge conflict resolution, and distributed development workflows.

\section{Conclusion}

In conclusion, our introduction of \benchmarkname~marks a significant step forward in the rigorous and realistic evaluation of AI coding agents. By adhering to three core principles—diverse, real-world task selection; challenging, multi-file code changes; and strict contamination prevention—we have created a benchmark that more accurately reflects the complexity of professional software engineering. Our findings, which show top-tier models like Opus 4.1 and GPT-5 achieving a 23\% success rate on \benchmarkname~compared to over 70\% on benchmarks like SWE-Bench Verified, highlight a critical gap between current agent capabilities and the demands of real-world development. This new baseline not only provides a more accurate measure of progress but also offers crucial insights into the specific limitations that must be addressed to advance the field. \benchmarkname serves as a robust, contamination-resistant testbed that can help guide future research toward developing truly autonomous and capable software engineering agents.

\newpage

% \section*{Acknowledgments}

% We would like to thank the contributors for their hard work on their dataset. Some of them are: Fernando Carabedo, Donnahue George Jr, Elías Muñiz. We are also deeply appreciative of the early-stage startups that partnered with us to provide proprietary commercial codebases, enabling a more realistic evaluation of AI agents in enterprise settings. Finally, we acknowledge the open-source communities behind the GPL-licensed repositories for their foundational work in software engineering, which inspired this benchmark. This research would not have been possible without these collective efforts.